% !TEX root = ../main.tex
%% Préparer l'épreuve : format de l'écrit, grille de rédaction, ECE et fiches expérimentales

\editionPage{epreuve-guide}{Format de l'épreuve et conseils}{Préparer la partie écrite}{Session 2027 · Note de service du 11 septembre 2026 (NOR MENE2622644N)}
\begin{blocEdition}{epreuve-guide}
\fontsize{10}{13.5}\selectfont
\editionEtape{L'épreuve en bref}
\begin{tabularx}{\linewidth}{@{}p{38mm}Y@{}}\toprule
Partie écrite & 3 h 30 ; trois exercices indépendants, à traiter tous ; note sur 20 points, dont deux consacrés à la
maîtrise de la langue. Le sujet précise si la calculatrice est autorisée (en mode examen).\\[1mm]
Partie pratique (ECE) & 1 h ; une situation d'évaluation tirée au sort ; note sur 20 points.\\[1mm]
Note de la spécialité & $0{,}8\times$ note de l'écrit $+\ 0{,}2\times$ note de l'ECE. Les candidats individuels, du Cned et des
établissements hors contrat, dispensés d'ECE, sont notés sur l'écrit seul.\\[1mm]
Programme & Tout le programme de terminale. Les notions de première non approfondies en terminale doivent être connues,
mais ne peuvent pas constituer un ressort essentiel du sujet.\\[1mm]
Oral de contrôle & 20 minutes de préparation, 20 minutes d'oral ; deux questions portant sur deux domaines différents.\\\bottomrule
\end{tabularx}
\editionEtape{Ce que demandent les sujets}
Les exercices partent de situations contextualisées et s'appuient sur des documents : textes, schémas, graphiques,
tableaux de mesures, programmes en Python. Ils accordent une place importante à la \textbf{modélisation} et aux questions avec
\textbf{prise d'initiative}, signalées dans l'énoncé : aucune démarche n'y est imposée et la démarche, même inaboutie, est
évaluée. Ils comportent aussi des questions sur les \textbf{aspects expérimentaux} (protocole, matériel, incertitudes) et sur
les \textbf{capacités numériques} du programme (lire, compléter ou interpréter un programme).
\editionEtape{Gérer les trois heures trente}
\begin{questions}
\item Lire tout le sujet pendant une dizaine de minutes : repérer les données, les documents, les barèmes et les questions
indépendantes.
\item Répartir le temps selon le barème, environ 9 à 10 minutes par point : près d'1 h 25 pour un exercice sur 9 points,
55 minutes pour 6 points, 50 minutes pour 5 points ; garder une dizaine de minutes pour relire.
\item Ne pas rester bloqué : une question « Montrer que » donne le résultat à utiliser pour la suite.
\end{questions}
\end{blocEdition}

\editionPage{epreuve-guide-2}{Rédiger une réponse}{Consignes, calculs et prise d'initiative}{Comprendre ce qui est demandé · Présenter ce qui est évalué}
\begin{blocEdition}{epreuve-guide-2}
\fontsize{10}{13.5}\selectfont
\editionEtape{Comprendre les verbes de consigne}
\begin{tabularx}{\linewidth}{@{}p{38mm}Y@{}}\toprule
Calculer & obtenir une valeur numérique, avec son unité et un nombre de chiffres significatifs adapté. \\
Déterminer & obtenir un résultat par la méthode de son choix (calcul, lecture graphique), en l'explicitant. \\
Établir, montrer que & mener une démonstration complète ; le résultat est connu, seul le raisonnement est évalué. \\
Exprimer & donner une expression littérale, en fonction des grandeurs indiquées. \\
Justifier & s'appuyer sur une loi, une propriété, une donnée ou un document. \\
Estimer, évaluer & donner un ordre de grandeur, avec des hypothèses simplificatrices. \\
Interpréter, commenter & confronter un résultat à un modèle, à une référence ou au contexte, puis conclure. \\\bottomrule
\end{tabularx}
\editionEtape{Rédiger un calcul en cinq temps}
\begin{questions}
\item Annoncer la loi ou la relation utilisée : « D'après la loi de Beer-Lambert\ldots ».
\item Établir l'expression littérale de la grandeur cherchée, avant toute application numérique.
\item Remplacer par les valeurs, converties dans des unités cohérentes (le plus souvent celles du Système international).
\item Écrire le résultat avec son unité et un nombre adapté de chiffres significatifs, puis le mettre en évidence.
\item Conclure par une phrase qui répond à la question posée.
\end{questions}
\editionEtape{Réussir une question avec prise d'initiative}
Faire un schéma de la situation ; lister les données utiles des documents ; nommer la grandeur cherchée ; choisir les lois
qui la relient aux données ; mener le calcul ; vérifier l'ordre de grandeur et conclure. Une démarche cohérente, même
incomplète, rapporte des points : il faut donc l'écrire.
\end{blocEdition}

\editionPage{epreuve-guide-3}{Grille de rédaction}{Les deux points de maîtrise de la langue}{Relire sa copie avant de la rendre}
\begin{blocEdition}{epreuve-guide-3}
\fontsize{10}{13.5}\selectfont
\editionEtape{Grille de relecture de la maîtrise de la langue}
La note de service consacre deux points à la maîtrise des normes orthographiques et syntaxiques, à la capacité à formuler
un raisonnement et à l'usage d'un vocabulaire juste et adapté. La grille ci-dessous, établie par OptiTAB d'après ces
attendus, permet de relire sa copie ; la grille officielle est annexée à la note de service.\par\smallskip
\begin{tabularx}{\linewidth}{@{}p{42mm}Y@{}}\toprule
\textbf{Attendu} & \textbf{Ce que le correcteur observe} \\\midrule
Normes orthographiques et syntaxiques & Phrases complètes et correctement construites ; accords ; ponctuation ; abréviations
évitées. \\[1mm]
Formuler un raisonnement & Enchaînement logique explicite (« d'après », « or », « donc ») ; hypothèses énoncées ; conclusion
qui répond à la question. \\[1mm]
Vocabulaire juste et adapté & Termes scientifiques précis : espèce et entité, transformation et réaction, grandeur et unité,
masse et quantité de matière, vitesse et accélération. \\[1mm]
Présentation & Questions numérotées ; résultats mis en évidence ; schémas légendés ; graphiques avec axes, grandeurs
et unités. \\\bottomrule
\end{tabularx}
\editionEtape{Les erreurs les plus coûteuses}
Un résultat sans unité ou avec trop de chiffres ; un oubli de conversion (mL en L, \unit{\degreeCelsius} en K, min en s, km en m) ;
une confusion entre concentration en masse et concentration en quantité de matière ; une méthode graphique non
expliquée (tangentes, 63~\%, demi-équivalence) ; un résultat aberrant non commenté.
\editionRetenir{Une copie claire se lit vite : une idée par phrase, une loi citée avant chaque calcul, un résultat encadré avec
son unité, une conclusion qui reprend les mots de la question.}
\end{blocEdition}

\editionPage{ece-guide}{Préparer l'ECE}{Évaluation des compétences expérimentales}{Une heure au laboratoire · 20~\% de la note de la spécialité}
\begin{blocEdition}{ece-guide}
\fontsize{10}{13.5}\selectfont
\editionEtape{Le déroulement}
L'évaluation dure une heure, dans une salle de travaux pratiques. Le candidat tire au sort sa situation d'évaluation parmi
au moins deux situations à dominante physique et deux à dominante chimie, et découvre le sujet en entrant dans la salle.
Deux professeurs examinateurs sont présents ; chacun évalue au plus quatre candidats et n'évalue pas ses propres élèves de
l'année. Le sujet comporte des \textbf{appels} : le candidat présente alors à l'examinateur un protocole, une mesure ou un
résultat. En cas de blocage, il peut demander une aide, ce que l'évaluation prend en compte.
\editionEtape{Les compétences évaluées}
\begin{tabularx}{\linewidth}{@{}p{38mm}Y@{}}\toprule
S'approprier & Comprendre la problématique, extraire l'information utile des documents, du matériel et du sujet. \\[1mm]
Analyser, raisonner & Proposer ou compléter un protocole, choisir le matériel, prévoir l'allure d'un résultat. \\[1mm]
Réaliser & Mettre en œuvre le protocole avec soin et en sécurité ; régler les appareils ; effectuer et noter les mesures. \\[1mm]
Valider & Exploiter les mesures, évaluer les incertitudes, comparer à une référence, identifier les sources d'erreur. \\[1mm]
Communiquer & Présenter clairement une démarche, un tableau de mesures, un graphique et une conclusion, à l'écrit et à
l'oral. \\\bottomrule
\end{tabularx}
\editionEtape{Le jour de l'épreuve}
\begin{questions}
\item Lire tout le sujet avant de toucher au matériel ; repérer les appels et le temps qu'ils demandent.
\item Vérifier le matériel et les règles de sécurité : blouse, lunettes, gants si nécessaire, produits identifiés.
\item Noter chaque mesure aussitôt, avec son unité ; pour un titrage, faire un titrage rapide puis un titrage précis.
\item Préparer à l'avance ce que l'on dira à l'appel : l'objectif, la méthode, le résultat attendu.
\item Évaluer l'incertitude, comparer à la valeur attendue et discuter les écarts.
\item Ranger le poste et éliminer les déchets dans les récipients prévus.
\end{questions}
\end{blocEdition}

\editionPage{ece-guide-2}{Capacités expérimentales}{Ce qu'il faut savoir faire au laboratoire}{Programme du cycle terminal · Renvois aux chapitres du livre}
\begin{blocEdition}{ece-guide-2}
\fontsize{10}{13.5}\selectfont
Trois capacités sont communes à tous les thèmes : respecter les règles de sécurité ; mettre en œuvre un dispositif
d'acquisition et de traitement de données (microcontrôleur, interface d'acquisition, tableur, Python) ; utiliser un
logiciel de simulation.
\editionEtape{Constitution et transformations de la matière}
\begin{tabularx}{\linewidth}{@{}Yr@{}}\toprule
Préparer une solution par dissolution ou par dilution en choisissant le matériel adapté & ch. 1, 3 \\
Réaliser des mesures d'absorbance, de pH, de conductivité ; tracer une courbe d'étalonnage & ch. 1, 2 \\
Réaliser un spectre UV-visible ; exploiter un spectre infrarouge & ch. 2, 11 \\
Mettre en œuvre le protocole d'un titrage (suivi pH-métrique, conductimétrique, colorimétrique) & ch. 3 \\
Suivre l'évolution temporelle d'une transformation & ch. 4 \\
Réaliser une pile et un circuit comportant un électrolyseur & ch. 8, 10 \\
Chauffage à reflux, distillation, extraction liquide-liquide, filtration, lavage, séchage & ch. 12 \\
Chromatographie sur couche mince ; température de changement d'état & ch. 12 \\\bottomrule
\end{tabularx}
\editionEtape{Mouvement, énergie, ondes et signaux}
\begin{tabularx}{\linewidth}{@{}Yr@{}}\toprule
Collecter des données sur un mouvement (vidéo, chronophotographie) et les exploiter & ch. 13 à 15 \\
Étudier la poussée d'Archimède ; mesurer une pression et une vitesse d'écoulement & ch. 17 \\
Mesurer une pression, un volume, une température ; tester la loi des gaz parfaits & ch. 18 \\
Réaliser un bilan énergétique et suivre l'évolution d'une température & ch. 19, 20 \\
Mesurer un niveau d'intensité sonore & ch. 21 \\
Mesurer la période, la longueur d'onde, la célérité d'une onde ; diffraction et interférences & ch. 22 \\
Mettre en évidence l'effet Doppler et mesurer une vitesse & ch. 23 \\
Estimer une distance focale ; réaliser une lunette astronomique & ch. 24 \\
Utiliser un multimètre ; réaliser un montage d'après un schéma normalisé ; circuit RC & ch. 25, 26 \\
Évaluer une incertitude-type, comparer à une valeur de référence & ch. 27 \\\bottomrule
\end{tabularx}
\editionRetenir{À l'ECE, une mesure n'est jamais seule : elle s'accompagne de son unité, de son incertitude et d'une phrase qui
dit si elle est compatible avec la valeur attendue.}
\end{blocEdition}

\editionPage{fiches-experimentales}{Fiche 1 · La verrerie}{Fiches expérimentales : verrerie, montages, mesures}{Choisir le bon instrument · Préparer une solution}
\begin{blocEdition}{fiches-experimentales}
\fontsize{10}{13.5}\selectfont
\label{figedition-fiches-1}\schema{annexe-verrerie}
\editionEtape{Précision de la verrerie}
{\small\begin{tabularx}{\linewidth}{@{}lYl@{}}\toprule
\textbf{Verrerie} & \textbf{Usage} & \textbf{Tolérance usuelle (classe A)} \\\midrule
Fiole jaugée de \qty{100,0}{mL} & préparer un volume précis de solution & $\pm\qty{0,10}{mL}$ \\
Pipette jaugée de \qty{10,00}{mL} & prélever un volume précis & $\pm\qty{0,02}{mL}$ \\
Burette graduée de \qty{25}{mL} & verser un volume variable (titrage) & $\pm\qty{0,03}{mL}$, lecture à \qty{0,05}{mL} \\
Éprouvette graduée de \qty{100}{mL} & mesurer un volume approximatif & $\pm\qty{0,5}{mL}$ \\
Bécher, erlenmeyer & contenir, mélanger, chauffer & graduations indicatives \\\bottomrule
\end{tabularx}}\par\smallskip
L'incertitude-type associée à une tolérance $\pm a$ vaut $a/\sqrt3$ : \qty{0,058}{mL} pour la fiole de \qty{100,0}{mL}.
\editionEtape{Préparer une solution par dilution}
\begin{questions}
\item Calculer le volume de solution mère à prélever : $V_\text{mère}=\dfrac{c_\text{fille}\,V_\text{fille}}{c_\text{mère}}$ ; le facteur de dilution vaut
$F=\dfrac{c_\text{mère}}{c_\text{fille}}=\dfrac{V_\text{fille}}{V_\text{mère}}$.
\item Verser un peu de solution mère dans un bécher ; la prélever avec une pipette jaugée rincée avec cette solution.
\item La verser dans une fiole jaugée, compléter avec de l'eau distillée aux trois quarts, agiter.
\item Ajuster au trait de jauge (bas du ménisque, à hauteur des yeux), boucher et homogénéiser.
\end{questions}
\editionEtape{Préparer une solution par dissolution}
Peser la masse $m=c\,V\,M$ de soluté dans une coupelle sur une balance tarée ; l'introduire dans la fiole jaugée avec un
entonnoir, rincer coupelle et entonnoir à l'eau distillée, dissoudre, puis compléter au trait de jauge et homogénéiser.
\end{blocEdition}

\editionPage{fiches-experimentales-2}{Fiche 2 · Synthèse organique}{Fiches expérimentales : verrerie, montages, mesures}{Chauffer, séparer, purifier, identifier}
\begin{blocEdition}{fiches-experimentales-2}
\fontsize{10}{13.5}\selectfont
\label{figedition-fiches-2}\schema{12-montage}
\editionEtape{Les étapes d'une synthèse}
\begin{tabularx}{\linewidth}{@{}p{30mm}Y@{}}\toprule
Transformation & Chauffage à reflux : la température accélère la réaction, le réfrigérant condense les vapeurs, sans
perte de matière. Pierre ponce pour régulariser l'ébullition ; catalyseur éventuel. \\[1mm]
Séparation & Extraction liquide-liquide en ampoule à décanter (solvant non miscible à l'eau, dans lequel le produit est
très soluble) ; filtration simple ou sous pression réduite (Büchner) pour un solide. \\[1mm]
Purification & Lavage de la phase organique, séchage par un desséchant anhydre ; recristallisation d'un solide ;
distillation fractionnée d'un liquide. \\[1mm]
Identification & Température de fusion au banc Kofler ; chromatographie sur couche mince ; spectroscopie infrarouge ;
indice de réfraction. \\\bottomrule
\end{tabularx}
\editionEtape{Réaliser une chromatographie sur couche mince}
Déposer, sur la ligne de dépôt tracée au crayon, le produit obtenu et les références ; placer la plaque dans la cuve
saturée en vapeurs d'éluant, la ligne de dépôt au-dessus du niveau de l'éluant ; sortir la plaque avant que le front
n'atteigne le haut, marquer le front, sécher, révéler (lampe UV, diiode). Deux espèces identiques migrent à la même
hauteur ; un produit pur ne donne qu'une tache.
\editionEtape{Rendement et sécurité}
Rendement : $\eta=\dfrac{n_\text{obtenue}}{n_\text{max}}$, où $n_\text{max}$ se calcule à partir du réactif limitant. Porter blouse, lunettes et
gants ; manipuler les espèces volatiles ou corrosives sous la hotte ; consulter les pictogrammes de danger ; éliminer
les déchets dans les récipients prévus.
\end{blocEdition}

\editionPage{fiches-experimentales-3}{Fiche 3 · Titrages et spectrophotométrie}{Fiches expérimentales : verrerie, montages, mesures}{Doser une espèce chimique en solution}
\begin{blocEdition}{fiches-experimentales-3}
\fontsize{10}{13.5}\selectfont
\label{figedition-fiches-3}\schema{03-montage-titrage}
\editionEtape{Mener un titrage}
Étalonner le pH-mètre avec deux solutions tampons (ou rincer la cellule du conductimètre) ; prélever la solution titrée à
la pipette jaugée ; remplir la burette, chasser la bulle d'air, ajuster le zéro. Verser la solution titrante millilitre
par millilitre, puis goutte à goutte au voisinage de l'équivalence, en notant le pH ou la conductivité après chaque ajout.
\editionEtape{Repérer l'équivalence}
\begin{tabularx}{\linewidth}{@{}p{38mm}Y@{}}\toprule
Suivi pH-métrique & Méthode des tangentes parallèles, ou maximum de la dérivée $\dfrac{\mathrm{d}\pH}{\mathrm{d}V}$. \\[1mm]
Suivi conductimétrique & Intersection des deux portions de droite de la courbe $\sigma=f(V)$. \\[1mm]
Indicateur coloré & Changement de couleur persistant ; la zone de virage doit être contenue dans le saut de pH. \\\bottomrule
\end{tabularx}\par\smallskip
À l'équivalence, les réactifs ont été introduits dans les proportions stœchiométriques : pour $a\,\mathrm{A}+b\,\mathrm{B}\to\ldots$,
$\dfrac{n_\mathrm{A}}{a}=\dfrac{n_{\mathrm{B},\mathrm{E}}}{b}$.
\editionEtape{Doser par spectrophotométrie}
Faire le blanc avec le solvant ; se placer à la longueur d'onde du maximum d'absorption ; mesurer l'absorbance de
solutions étalons et tracer la droite d'étalonnage $A=k\,c$ (loi de Beer-Lambert, valable pour des solutions diluées) ;
mesurer l'absorbance de la solution inconnue, diluée si nécessaire pour rester dans le domaine d'étalonnage, puis en
déduire sa concentration.
\end{blocEdition}

\editionPage{fiches-experimentales-4}{Fiche 4 · Mesures en physique}{Fiches expérimentales : verrerie, montages, mesures}{Acquérir, traiter, évaluer les incertitudes}
\begin{blocEdition}{fiches-experimentales-4}
\fontsize{10}{13.5}\selectfont
\editionEtape{Acquérir un signal}
\begin{tabularx}{\linewidth}{@{}p{38mm}Y@{}}\toprule
Oscilloscope, interface & Régler la base de temps pour voir quelques périodes ; mesurer une période sur plusieurs
périodes ; un retard entre deux signaux se mesure entre deux points homologues. \\[1mm]
Microcontrôleur & Commander un signal, lire un capteur, horodater les mesures ; vérifier la fréquence d'échantillonnage. \\[1mm]
Pointage vidéo & Placer une échelle de longueur dans le plan du mouvement ; repérer l'origine des dates ; exporter les
coordonnées vers un tableur ou Python. \\[1mm]
Multimètre & Brancher un voltmètre en dérivation, un ampèremètre en série ; commencer par le plus grand calibre. \\[1mm]
Sonomètre, luxmètre & Mesurer à distance fixe de la source, sans obstacle, en notant le bruit de fond. \\\bottomrule
\end{tabularx}
\editionEtape{Optique}
Distance focale d'une lentille convergente : former sur un écran l'image nette d'un objet éloigné, ou utiliser la relation
de conjugaison $\dfrac{1}{\overline{OA'}}-\dfrac{1}{\overline{OA}}=\dfrac{1}{f'}$ pour plusieurs positions. Lunette afocale : placer le foyer objet de
l'oculaire au foyer image de l'objectif.
\editionEtape{Évaluer une incertitude}
\begin{tabularx}{\linewidth}{@{}p{38mm}Y@{}}\toprule
Série de $N$ mesures (type A) & $u(\bar m)=\dfrac{s}{\sqrt N}$, $s$ étant l'écart-type expérimental. \\[1mm]
Mesure unique (type B) & $u=\dfrac{a}{\sqrt3}$, $a$ étant la demi-étendue (demi-graduation, tolérance du constructeur). \\[1mm]
Grandeur calculée & Expression fournie, ou simulation de Monte-Carlo en Python. \\[1mm]
Comparaison & $z=\dfrac{|m_\text{mes}-m_\text{ref}|}{u(m)}$ ; compatible si $z<2$. \\\bottomrule
\end{tabularx}
\editionEtape{Python au laboratoire}
\code{mean} et \code{stdev} (module \code{statistics}) donnent moyenne et écart-type ; \code{np.polyfit(x, y, 1)} réalise une
régression linéaire ; \code{plt.plot}, \code{plt.hist} et \code{plt.show} tracent des courbes et des histogrammes ; la méthode
d'Euler résout une équation différentielle pas à pas (chapitres 14, 20 et 26).
\end{blocEdition}
